\documentclass[11pt, a4paper]{article}
% =============================================================
%  preamble.tex — Reusable preamble for Neural Networks course
%  Usage: % =============================================================
%  preamble.tex — Reusable preamble for Neural Networks course
%  Usage: % =============================================================
%  preamble.tex — Reusable preamble for Neural Networks course
%  Usage: \input{preamble} at the top of your document,
%         after \documentclass but before \begin{document}.
% =============================================================

% --- Encoding & language ---
\usepackage[T1]{fontenc}
\usepackage[utf8]{inputenc}
\usepackage[english]{babel}         % swap for [french] if needed

% --- Page geometry ---
\usepackage{geometry}
\geometry{top=2.5cm, bottom=2.5cm, left=2.5cm, right=2.5cm}

% --- Math ---
\usepackage{amsmath, amssymb}

% --- Tables ---
\usepackage{booktabs}               % \toprule, \midrule, \bottomrule
\usepackage{array}                  % column type extensions
\usepackage{tabularx}               % auto-width columns
\usepackage{multirow}               % multi-row cells

% --- Colors ---
\usepackage{xcolor}

\definecolor{mainblue}{RGB}{30, 80, 160}
\definecolor{lightblue}{RGB}{220, 235, 255}
\definecolor{accentgray}{RGB}{245, 246, 248}
\definecolor{darkgray}{RGB}{80, 80, 80}
\definecolor{solutiongreen}{RGB}{20, 120, 60}
\definecolor{lightgreen}{RGB}{220, 245, 230}
\definecolor{warnorange}{RGB}{210, 120, 20}
\definecolor{lightorange}{RGB}{255, 243, 220}

% --- Lists ---
\usepackage{enumitem}

% --- TikZ & pgfplots ---
\usepackage{tikz}
\usepackage{pgfplots}
\pgfplotsset{compat=1.18}
\usetikzlibrary{positioning, arrows.meta, calc, patterns}

% --- Boxes: tcolorbox ---
\usepackage{tcolorbox}
\tcbuselibrary{skins, breakable}

% Generic exercise box (blue border, gray background)
\tcbset{
  exercisebox/.style={
    enhanced,
    colback=accentgray,
    colframe=mainblue,
    boxrule=0.8pt,
    arc=4pt,
    left=10pt, right=10pt, top=8pt, bottom=8pt,
    fonttitle=\bfseries\color{mainblue},
  }
}

% Solution box (green, used in answer sheets)
\tcbset{
  solutionbox/.style={
    enhanced,
    colback=lightgreen,
    colframe=solutiongreen,
    boxrule=0.8pt,
    arc=4pt,
    left=10pt, right=10pt, top=8pt, bottom=8pt,
    fonttitle=\bfseries\color{solutiongreen},
  }
}

% Answer highlight box (orange, for final numeric answers)
\tcbset{
  answerbox/.style={
    enhanced,
    colback=lightorange,
    colframe=warnorange,
    boxrule=0.7pt,
    arc=4pt,
    left=8pt, right=8pt, top=5pt, bottom=5pt,
    fonttitle=\bfseries\color{warnorange},
  }
}

% Step box (light blue border, gray background, for intermediate steps)
\tcbset{
  stepbox/.style={
    enhanced,
    colback=accentgray,
    colframe=mainblue!40,
    boxrule=0.5pt,
    arc=3pt,
    left=8pt, right=8pt, top=6pt, bottom=6pt,
  }
}

% --- Inline note / warning box (mdframed) ---
\usepackage{mdframed}
\newmdenv[
  linecolor=warnorange,
  backgroundcolor=orange!5,
  roundcorner=4pt,
  innerleftmargin=8pt,
  innerrightmargin=8pt,
  innertopmargin=6pt,
  innerbottommargin=6pt,
  linewidth=0.7pt
]{notebox}

% --- Section formatting ---
\usepackage{titlesec}

% Section style: "Exercise N : Title" with a blue rule underneath
% Change the prefix via \renewcommand{\sectionprefix}{...} before \section
\newcommand{\sectionprefix}{Exercise}
\titleformat{\section}[block]
  {\large\bfseries\color{mainblue}}
  {\sectionprefix~\thesection\,:}{0.6em}{}
  [\vspace{-0.3em}\textcolor{mainblue}{\rule{\linewidth}{0.8pt}}]
\titlespacing{\section}{0pt}{1.6em}{0.8em}

% --- Header / Footer ---
\usepackage{fancyhdr}
\pagestyle{fancy}
\fancyhf{}
\renewcommand{\headrulewidth}{0.4pt}

% Customise these three commands in your document after \input{preamble}:
%   \renewcommand{\doccoursename}{My Course}
%   \renewcommand{\docsessionname}{Session N — Topic}
\newcommand{\doccoursename}{Neural Networks for Engineers}
\newcommand{\docsessionname}{Session}

\fancyhead[L]{\small\textcolor{darkgray}{\doccoursename}}
\fancyhead[R]{\small\textcolor{darkgray}{\docsessionname}}
\fancyfoot[C]{\small\textcolor{darkgray}{\thepage}}

% --- Title block macro ---
% Usage: \maketitleblock{Title line}{Subtitle / metadata line}{background color}
% Examples:
%   \maketitleblock{Session 4 — Exercises}{Mixed level | Neural Networks}{mainblue}
%   \maketitleblock{Session 4 — Solutions}{Perceptrons and MLPs | Neural Networks}{solutiongreen}
\newcommand{\maketitleblock}[3]{%
  \begin{tcolorbox}[
    enhanced,
    colback=#3,
    colframe=#3,
    arc=6pt, boxrule=0pt,
    left=16pt, right=16pt, top=12pt, bottom=12pt
  ]
    {\large\bfseries\color{white} #1}\\[4pt]
    {\small\color{white!85!black} #2}
  \end{tcolorbox}
  \vspace{1em}
}

% --- Convenience macros ---

% Blank answer field (inline underline for fill-in tables)
\newcommand{\acell}{\rule{1.8cm}{0.4pt}}

% Step-by-step computation display: \step{label}{math}
% Example: \step{Hidden layer}{z^{(1)} = Wx + b = \ldots}
\newcommand{\step}[2]{%
  \noindent\textbf{#1:}\quad $#2$\par\vspace{0.3em}
}

% Neuron style for TikZ diagrams (reuse with \node[neuron])
% Colors are overridable per layer — see examples below.
%
% Input neuron  : fill=lightblue,  draw=mainblue
% Hidden neuron : fill=green!15,   draw=green!50!black
% Output neuron : fill=orange!15,  draw=orange!70!black
%
% Example architecture diagram:
%
%   \begin{tikzpicture}[
%     neuron/.style={circle, minimum size=1cm, thick, font=\small\bfseries},
%     arr/.style={->, >=Stealth, thick, gray!70}
%   ]
%     \node[neuron, fill=lightblue,  draw=mainblue]          (x1) at (0, 1)  {$x_1$};
%     \node[neuron, fill=green!15,   draw=green!50!black]    (h1) at (3, 1)  {$h_1$};
%     \node[neuron, fill=orange!15,  draw=orange!70!black]   (y)  at (6, 1)  {$\hat{y}$};
%     \draw[arr] (x1) -- (h1);
%     \draw[arr] (h1) -- (y);
%   \end{tikzpicture}
 at the top of your document,
%         after \documentclass but before \begin{document}.
% =============================================================

% --- Encoding & language ---
\usepackage[T1]{fontenc}
\usepackage[utf8]{inputenc}
\usepackage[english]{babel}         % swap for [french] if needed

% --- Page geometry ---
\usepackage{geometry}
\geometry{top=2.5cm, bottom=2.5cm, left=2.5cm, right=2.5cm}

% --- Math ---
\usepackage{amsmath, amssymb}

% --- Tables ---
\usepackage{booktabs}               % \toprule, \midrule, \bottomrule
\usepackage{array}                  % column type extensions
\usepackage{tabularx}               % auto-width columns
\usepackage{multirow}               % multi-row cells

% --- Colors ---
\usepackage{xcolor}

\definecolor{mainblue}{RGB}{30, 80, 160}
\definecolor{lightblue}{RGB}{220, 235, 255}
\definecolor{accentgray}{RGB}{245, 246, 248}
\definecolor{darkgray}{RGB}{80, 80, 80}
\definecolor{solutiongreen}{RGB}{20, 120, 60}
\definecolor{lightgreen}{RGB}{220, 245, 230}
\definecolor{warnorange}{RGB}{210, 120, 20}
\definecolor{lightorange}{RGB}{255, 243, 220}

% --- Lists ---
\usepackage{enumitem}

% --- TikZ & pgfplots ---
\usepackage{tikz}
\usepackage{pgfplots}
\pgfplotsset{compat=1.18}
\usetikzlibrary{positioning, arrows.meta, calc, patterns}

% --- Boxes: tcolorbox ---
\usepackage{tcolorbox}
\tcbuselibrary{skins, breakable}

% Generic exercise box (blue border, gray background)
\tcbset{
  exercisebox/.style={
    enhanced,
    colback=accentgray,
    colframe=mainblue,
    boxrule=0.8pt,
    arc=4pt,
    left=10pt, right=10pt, top=8pt, bottom=8pt,
    fonttitle=\bfseries\color{mainblue},
  }
}

% Solution box (green, used in answer sheets)
\tcbset{
  solutionbox/.style={
    enhanced,
    colback=lightgreen,
    colframe=solutiongreen,
    boxrule=0.8pt,
    arc=4pt,
    left=10pt, right=10pt, top=8pt, bottom=8pt,
    fonttitle=\bfseries\color{solutiongreen},
  }
}

% Answer highlight box (orange, for final numeric answers)
\tcbset{
  answerbox/.style={
    enhanced,
    colback=lightorange,
    colframe=warnorange,
    boxrule=0.7pt,
    arc=4pt,
    left=8pt, right=8pt, top=5pt, bottom=5pt,
    fonttitle=\bfseries\color{warnorange},
  }
}

% Step box (light blue border, gray background, for intermediate steps)
\tcbset{
  stepbox/.style={
    enhanced,
    colback=accentgray,
    colframe=mainblue!40,
    boxrule=0.5pt,
    arc=3pt,
    left=8pt, right=8pt, top=6pt, bottom=6pt,
  }
}

% --- Inline note / warning box (mdframed) ---
\usepackage{mdframed}
\newmdenv[
  linecolor=warnorange,
  backgroundcolor=orange!5,
  roundcorner=4pt,
  innerleftmargin=8pt,
  innerrightmargin=8pt,
  innertopmargin=6pt,
  innerbottommargin=6pt,
  linewidth=0.7pt
]{notebox}

% --- Section formatting ---
\usepackage{titlesec}

% Section style: "Exercise N : Title" with a blue rule underneath
% Change the prefix via \renewcommand{\sectionprefix}{...} before \section
\newcommand{\sectionprefix}{Exercise}
\titleformat{\section}[block]
  {\large\bfseries\color{mainblue}}
  {\sectionprefix~\thesection\,:}{0.6em}{}
  [\vspace{-0.3em}\textcolor{mainblue}{\rule{\linewidth}{0.8pt}}]
\titlespacing{\section}{0pt}{1.6em}{0.8em}

% --- Header / Footer ---
\usepackage{fancyhdr}
\pagestyle{fancy}
\fancyhf{}
\renewcommand{\headrulewidth}{0.4pt}

% Customise these three commands in your document after % =============================================================
%  preamble.tex — Reusable preamble for Neural Networks course
%  Usage: \input{preamble} at the top of your document,
%         after \documentclass but before \begin{document}.
% =============================================================

% --- Encoding & language ---
\usepackage[T1]{fontenc}
\usepackage[utf8]{inputenc}
\usepackage[english]{babel}         % swap for [french] if needed

% --- Page geometry ---
\usepackage{geometry}
\geometry{top=2.5cm, bottom=2.5cm, left=2.5cm, right=2.5cm}

% --- Math ---
\usepackage{amsmath, amssymb}

% --- Tables ---
\usepackage{booktabs}               % \toprule, \midrule, \bottomrule
\usepackage{array}                  % column type extensions
\usepackage{tabularx}               % auto-width columns
\usepackage{multirow}               % multi-row cells

% --- Colors ---
\usepackage{xcolor}

\definecolor{mainblue}{RGB}{30, 80, 160}
\definecolor{lightblue}{RGB}{220, 235, 255}
\definecolor{accentgray}{RGB}{245, 246, 248}
\definecolor{darkgray}{RGB}{80, 80, 80}
\definecolor{solutiongreen}{RGB}{20, 120, 60}
\definecolor{lightgreen}{RGB}{220, 245, 230}
\definecolor{warnorange}{RGB}{210, 120, 20}
\definecolor{lightorange}{RGB}{255, 243, 220}

% --- Lists ---
\usepackage{enumitem}

% --- TikZ & pgfplots ---
\usepackage{tikz}
\usepackage{pgfplots}
\pgfplotsset{compat=1.18}
\usetikzlibrary{positioning, arrows.meta, calc, patterns}

% --- Boxes: tcolorbox ---
\usepackage{tcolorbox}
\tcbuselibrary{skins, breakable}

% Generic exercise box (blue border, gray background)
\tcbset{
  exercisebox/.style={
    enhanced,
    colback=accentgray,
    colframe=mainblue,
    boxrule=0.8pt,
    arc=4pt,
    left=10pt, right=10pt, top=8pt, bottom=8pt,
    fonttitle=\bfseries\color{mainblue},
  }
}

% Solution box (green, used in answer sheets)
\tcbset{
  solutionbox/.style={
    enhanced,
    colback=lightgreen,
    colframe=solutiongreen,
    boxrule=0.8pt,
    arc=4pt,
    left=10pt, right=10pt, top=8pt, bottom=8pt,
    fonttitle=\bfseries\color{solutiongreen},
  }
}

% Answer highlight box (orange, for final numeric answers)
\tcbset{
  answerbox/.style={
    enhanced,
    colback=lightorange,
    colframe=warnorange,
    boxrule=0.7pt,
    arc=4pt,
    left=8pt, right=8pt, top=5pt, bottom=5pt,
    fonttitle=\bfseries\color{warnorange},
  }
}

% Step box (light blue border, gray background, for intermediate steps)
\tcbset{
  stepbox/.style={
    enhanced,
    colback=accentgray,
    colframe=mainblue!40,
    boxrule=0.5pt,
    arc=3pt,
    left=8pt, right=8pt, top=6pt, bottom=6pt,
  }
}

% --- Inline note / warning box (mdframed) ---
\usepackage{mdframed}
\newmdenv[
  linecolor=warnorange,
  backgroundcolor=orange!5,
  roundcorner=4pt,
  innerleftmargin=8pt,
  innerrightmargin=8pt,
  innertopmargin=6pt,
  innerbottommargin=6pt,
  linewidth=0.7pt
]{notebox}

% --- Section formatting ---
\usepackage{titlesec}

% Section style: "Exercise N : Title" with a blue rule underneath
% Change the prefix via \renewcommand{\sectionprefix}{...} before \section
\newcommand{\sectionprefix}{Exercise}
\titleformat{\section}[block]
  {\large\bfseries\color{mainblue}}
  {\sectionprefix~\thesection\,:}{0.6em}{}
  [\vspace{-0.3em}\textcolor{mainblue}{\rule{\linewidth}{0.8pt}}]
\titlespacing{\section}{0pt}{1.6em}{0.8em}

% --- Header / Footer ---
\usepackage{fancyhdr}
\pagestyle{fancy}
\fancyhf{}
\renewcommand{\headrulewidth}{0.4pt}

% Customise these three commands in your document after \input{preamble}:
%   \renewcommand{\doccoursename}{My Course}
%   \renewcommand{\docsessionname}{Session N — Topic}
\newcommand{\doccoursename}{Neural Networks for Engineers}
\newcommand{\docsessionname}{Session}

\fancyhead[L]{\small\textcolor{darkgray}{\doccoursename}}
\fancyhead[R]{\small\textcolor{darkgray}{\docsessionname}}
\fancyfoot[C]{\small\textcolor{darkgray}{\thepage}}

% --- Title block macro ---
% Usage: \maketitleblock{Title line}{Subtitle / metadata line}{background color}
% Examples:
%   \maketitleblock{Session 4 — Exercises}{Mixed level | Neural Networks}{mainblue}
%   \maketitleblock{Session 4 — Solutions}{Perceptrons and MLPs | Neural Networks}{solutiongreen}
\newcommand{\maketitleblock}[3]{%
  \begin{tcolorbox}[
    enhanced,
    colback=#3,
    colframe=#3,
    arc=6pt, boxrule=0pt,
    left=16pt, right=16pt, top=12pt, bottom=12pt
  ]
    {\large\bfseries\color{white} #1}\\[4pt]
    {\small\color{white!85!black} #2}
  \end{tcolorbox}
  \vspace{1em}
}

% --- Convenience macros ---

% Blank answer field (inline underline for fill-in tables)
\newcommand{\acell}{\rule{1.8cm}{0.4pt}}

% Step-by-step computation display: \step{label}{math}
% Example: \step{Hidden layer}{z^{(1)} = Wx + b = \ldots}
\newcommand{\step}[2]{%
  \noindent\textbf{#1:}\quad $#2$\par\vspace{0.3em}
}

% Neuron style for TikZ diagrams (reuse with \node[neuron])
% Colors are overridable per layer — see examples below.
%
% Input neuron  : fill=lightblue,  draw=mainblue
% Hidden neuron : fill=green!15,   draw=green!50!black
% Output neuron : fill=orange!15,  draw=orange!70!black
%
% Example architecture diagram:
%
%   \begin{tikzpicture}[
%     neuron/.style={circle, minimum size=1cm, thick, font=\small\bfseries},
%     arr/.style={->, >=Stealth, thick, gray!70}
%   ]
%     \node[neuron, fill=lightblue,  draw=mainblue]          (x1) at (0, 1)  {$x_1$};
%     \node[neuron, fill=green!15,   draw=green!50!black]    (h1) at (3, 1)  {$h_1$};
%     \node[neuron, fill=orange!15,  draw=orange!70!black]   (y)  at (6, 1)  {$\hat{y}$};
%     \draw[arr] (x1) -- (h1);
%     \draw[arr] (h1) -- (y);
%   \end{tikzpicture}
:
%   \renewcommand{\doccoursename}{My Course}
%   \renewcommand{\docsessionname}{Session N — Topic}
\newcommand{\doccoursename}{Neural Networks for Engineers}
\newcommand{\docsessionname}{Session}

\fancyhead[L]{\small\textcolor{darkgray}{\doccoursename}}
\fancyhead[R]{\small\textcolor{darkgray}{\docsessionname}}
\fancyfoot[C]{\small\textcolor{darkgray}{\thepage}}

% --- Title block macro ---
% Usage: \maketitleblock{Title line}{Subtitle / metadata line}{background color}
% Examples:
%   \maketitleblock{Session 4 — Exercises}{Mixed level | Neural Networks}{mainblue}
%   \maketitleblock{Session 4 — Solutions}{Perceptrons and MLPs | Neural Networks}{solutiongreen}
\newcommand{\maketitleblock}[3]{%
  \begin{tcolorbox}[
    enhanced,
    colback=#3,
    colframe=#3,
    arc=6pt, boxrule=0pt,
    left=16pt, right=16pt, top=12pt, bottom=12pt
  ]
    {\large\bfseries\color{white} #1}\\[4pt]
    {\small\color{white!85!black} #2}
  \end{tcolorbox}
  \vspace{1em}
}

% --- Convenience macros ---

% Blank answer field (inline underline for fill-in tables)
\newcommand{\acell}{\rule{1.8cm}{0.4pt}}

% Step-by-step computation display: \step{label}{math}
% Example: \step{Hidden layer}{z^{(1)} = Wx + b = \ldots}
\newcommand{\step}[2]{%
  \noindent\textbf{#1:}\quad $#2$\par\vspace{0.3em}
}

% Neuron style for TikZ diagrams (reuse with \node[neuron])
% Colors are overridable per layer — see examples below.
%
% Input neuron  : fill=lightblue,  draw=mainblue
% Hidden neuron : fill=green!15,   draw=green!50!black
% Output neuron : fill=orange!15,  draw=orange!70!black
%
% Example architecture diagram:
%
%   \begin{tikzpicture}[
%     neuron/.style={circle, minimum size=1cm, thick, font=\small\bfseries},
%     arr/.style={->, >=Stealth, thick, gray!70}
%   ]
%     \node[neuron, fill=lightblue,  draw=mainblue]          (x1) at (0, 1)  {$x_1$};
%     \node[neuron, fill=green!15,   draw=green!50!black]    (h1) at (3, 1)  {$h_1$};
%     \node[neuron, fill=orange!15,  draw=orange!70!black]   (y)  at (6, 1)  {$\hat{y}$};
%     \draw[arr] (x1) -- (h1);
%     \draw[arr] (h1) -- (y);
%   \end{tikzpicture}
 at the top of your document,
%         after \documentclass but before \begin{document}.
% =============================================================

% --- Encoding & language ---
\usepackage[T1]{fontenc}
\usepackage[utf8]{inputenc}
\usepackage[english]{babel}         % swap for [french] if needed

% --- Page geometry ---
\usepackage{geometry}
\geometry{top=2.5cm, bottom=2.5cm, left=2.5cm, right=2.5cm}

% --- Math ---
\usepackage{amsmath, amssymb}

% --- Tables ---
\usepackage{booktabs}               % \toprule, \midrule, \bottomrule
\usepackage{array}                  % column type extensions
\usepackage{tabularx}               % auto-width columns
\usepackage{multirow}               % multi-row cells

% --- Colors ---
\usepackage{xcolor}

\definecolor{mainblue}{RGB}{30, 80, 160}
\definecolor{lightblue}{RGB}{220, 235, 255}
\definecolor{accentgray}{RGB}{245, 246, 248}
\definecolor{darkgray}{RGB}{80, 80, 80}
\definecolor{solutiongreen}{RGB}{20, 120, 60}
\definecolor{lightgreen}{RGB}{220, 245, 230}
\definecolor{warnorange}{RGB}{210, 120, 20}
\definecolor{lightorange}{RGB}{255, 243, 220}

% --- Lists ---
\usepackage{enumitem}

% --- TikZ & pgfplots ---
\usepackage{tikz}
\usepackage{pgfplots}
\pgfplotsset{compat=1.18}
\usetikzlibrary{positioning, arrows.meta, calc, patterns}

% --- Boxes: tcolorbox ---
\usepackage{tcolorbox}
\tcbuselibrary{skins, breakable}

% Generic exercise box (blue border, gray background)
\tcbset{
  exercisebox/.style={
    enhanced,
    colback=accentgray,
    colframe=mainblue,
    boxrule=0.8pt,
    arc=4pt,
    left=10pt, right=10pt, top=8pt, bottom=8pt,
    fonttitle=\bfseries\color{mainblue},
  }
}

% Solution box (green, used in answer sheets)
\tcbset{
  solutionbox/.style={
    enhanced,
    colback=lightgreen,
    colframe=solutiongreen,
    boxrule=0.8pt,
    arc=4pt,
    left=10pt, right=10pt, top=8pt, bottom=8pt,
    fonttitle=\bfseries\color{solutiongreen},
  }
}

% Answer highlight box (orange, for final numeric answers)
\tcbset{
  answerbox/.style={
    enhanced,
    colback=lightorange,
    colframe=warnorange,
    boxrule=0.7pt,
    arc=4pt,
    left=8pt, right=8pt, top=5pt, bottom=5pt,
    fonttitle=\bfseries\color{warnorange},
  }
}

% Step box (light blue border, gray background, for intermediate steps)
\tcbset{
  stepbox/.style={
    enhanced,
    colback=accentgray,
    colframe=mainblue!40,
    boxrule=0.5pt,
    arc=3pt,
    left=8pt, right=8pt, top=6pt, bottom=6pt,
  }
}

% --- Inline note / warning box (mdframed) ---
\usepackage{mdframed}
\newmdenv[
  linecolor=warnorange,
  backgroundcolor=orange!5,
  roundcorner=4pt,
  innerleftmargin=8pt,
  innerrightmargin=8pt,
  innertopmargin=6pt,
  innerbottommargin=6pt,
  linewidth=0.7pt
]{notebox}

% --- Section formatting ---
\usepackage{titlesec}

% Section style: "Exercise N : Title" with a blue rule underneath
% Change the prefix via \renewcommand{\sectionprefix}{...} before \section
\newcommand{\sectionprefix}{Exercise}
\titleformat{\section}[block]
  {\large\bfseries\color{mainblue}}
  {\sectionprefix~\thesection\,:}{0.6em}{}
  [\vspace{-0.3em}\textcolor{mainblue}{\rule{\linewidth}{0.8pt}}]
\titlespacing{\section}{0pt}{1.6em}{0.8em}

% --- Header / Footer ---
\usepackage{fancyhdr}
\pagestyle{fancy}
\fancyhf{}
\renewcommand{\headrulewidth}{0.4pt}

% Customise these three commands in your document after % =============================================================
%  preamble.tex — Reusable preamble for Neural Networks course
%  Usage: % =============================================================
%  preamble.tex — Reusable preamble for Neural Networks course
%  Usage: \input{preamble} at the top of your document,
%         after \documentclass but before \begin{document}.
% =============================================================

% --- Encoding & language ---
\usepackage[T1]{fontenc}
\usepackage[utf8]{inputenc}
\usepackage[english]{babel}         % swap for [french] if needed

% --- Page geometry ---
\usepackage{geometry}
\geometry{top=2.5cm, bottom=2.5cm, left=2.5cm, right=2.5cm}

% --- Math ---
\usepackage{amsmath, amssymb}

% --- Tables ---
\usepackage{booktabs}               % \toprule, \midrule, \bottomrule
\usepackage{array}                  % column type extensions
\usepackage{tabularx}               % auto-width columns
\usepackage{multirow}               % multi-row cells

% --- Colors ---
\usepackage{xcolor}

\definecolor{mainblue}{RGB}{30, 80, 160}
\definecolor{lightblue}{RGB}{220, 235, 255}
\definecolor{accentgray}{RGB}{245, 246, 248}
\definecolor{darkgray}{RGB}{80, 80, 80}
\definecolor{solutiongreen}{RGB}{20, 120, 60}
\definecolor{lightgreen}{RGB}{220, 245, 230}
\definecolor{warnorange}{RGB}{210, 120, 20}
\definecolor{lightorange}{RGB}{255, 243, 220}

% --- Lists ---
\usepackage{enumitem}

% --- TikZ & pgfplots ---
\usepackage{tikz}
\usepackage{pgfplots}
\pgfplotsset{compat=1.18}
\usetikzlibrary{positioning, arrows.meta, calc, patterns}

% --- Boxes: tcolorbox ---
\usepackage{tcolorbox}
\tcbuselibrary{skins, breakable}

% Generic exercise box (blue border, gray background)
\tcbset{
  exercisebox/.style={
    enhanced,
    colback=accentgray,
    colframe=mainblue,
    boxrule=0.8pt,
    arc=4pt,
    left=10pt, right=10pt, top=8pt, bottom=8pt,
    fonttitle=\bfseries\color{mainblue},
  }
}

% Solution box (green, used in answer sheets)
\tcbset{
  solutionbox/.style={
    enhanced,
    colback=lightgreen,
    colframe=solutiongreen,
    boxrule=0.8pt,
    arc=4pt,
    left=10pt, right=10pt, top=8pt, bottom=8pt,
    fonttitle=\bfseries\color{solutiongreen},
  }
}

% Answer highlight box (orange, for final numeric answers)
\tcbset{
  answerbox/.style={
    enhanced,
    colback=lightorange,
    colframe=warnorange,
    boxrule=0.7pt,
    arc=4pt,
    left=8pt, right=8pt, top=5pt, bottom=5pt,
    fonttitle=\bfseries\color{warnorange},
  }
}

% Step box (light blue border, gray background, for intermediate steps)
\tcbset{
  stepbox/.style={
    enhanced,
    colback=accentgray,
    colframe=mainblue!40,
    boxrule=0.5pt,
    arc=3pt,
    left=8pt, right=8pt, top=6pt, bottom=6pt,
  }
}

% --- Inline note / warning box (mdframed) ---
\usepackage{mdframed}
\newmdenv[
  linecolor=warnorange,
  backgroundcolor=orange!5,
  roundcorner=4pt,
  innerleftmargin=8pt,
  innerrightmargin=8pt,
  innertopmargin=6pt,
  innerbottommargin=6pt,
  linewidth=0.7pt
]{notebox}

% --- Section formatting ---
\usepackage{titlesec}

% Section style: "Exercise N : Title" with a blue rule underneath
% Change the prefix via \renewcommand{\sectionprefix}{...} before \section
\newcommand{\sectionprefix}{Exercise}
\titleformat{\section}[block]
  {\large\bfseries\color{mainblue}}
  {\sectionprefix~\thesection\,:}{0.6em}{}
  [\vspace{-0.3em}\textcolor{mainblue}{\rule{\linewidth}{0.8pt}}]
\titlespacing{\section}{0pt}{1.6em}{0.8em}

% --- Header / Footer ---
\usepackage{fancyhdr}
\pagestyle{fancy}
\fancyhf{}
\renewcommand{\headrulewidth}{0.4pt}

% Customise these three commands in your document after \input{preamble}:
%   \renewcommand{\doccoursename}{My Course}
%   \renewcommand{\docsessionname}{Session N — Topic}
\newcommand{\doccoursename}{Neural Networks for Engineers}
\newcommand{\docsessionname}{Session}

\fancyhead[L]{\small\textcolor{darkgray}{\doccoursename}}
\fancyhead[R]{\small\textcolor{darkgray}{\docsessionname}}
\fancyfoot[C]{\small\textcolor{darkgray}{\thepage}}

% --- Title block macro ---
% Usage: \maketitleblock{Title line}{Subtitle / metadata line}{background color}
% Examples:
%   \maketitleblock{Session 4 — Exercises}{Mixed level | Neural Networks}{mainblue}
%   \maketitleblock{Session 4 — Solutions}{Perceptrons and MLPs | Neural Networks}{solutiongreen}
\newcommand{\maketitleblock}[3]{%
  \begin{tcolorbox}[
    enhanced,
    colback=#3,
    colframe=#3,
    arc=6pt, boxrule=0pt,
    left=16pt, right=16pt, top=12pt, bottom=12pt
  ]
    {\large\bfseries\color{white} #1}\\[4pt]
    {\small\color{white!85!black} #2}
  \end{tcolorbox}
  \vspace{1em}
}

% --- Convenience macros ---

% Blank answer field (inline underline for fill-in tables)
\newcommand{\acell}{\rule{1.8cm}{0.4pt}}

% Step-by-step computation display: \step{label}{math}
% Example: \step{Hidden layer}{z^{(1)} = Wx + b = \ldots}
\newcommand{\step}[2]{%
  \noindent\textbf{#1:}\quad $#2$\par\vspace{0.3em}
}

% Neuron style for TikZ diagrams (reuse with \node[neuron])
% Colors are overridable per layer — see examples below.
%
% Input neuron  : fill=lightblue,  draw=mainblue
% Hidden neuron : fill=green!15,   draw=green!50!black
% Output neuron : fill=orange!15,  draw=orange!70!black
%
% Example architecture diagram:
%
%   \begin{tikzpicture}[
%     neuron/.style={circle, minimum size=1cm, thick, font=\small\bfseries},
%     arr/.style={->, >=Stealth, thick, gray!70}
%   ]
%     \node[neuron, fill=lightblue,  draw=mainblue]          (x1) at (0, 1)  {$x_1$};
%     \node[neuron, fill=green!15,   draw=green!50!black]    (h1) at (3, 1)  {$h_1$};
%     \node[neuron, fill=orange!15,  draw=orange!70!black]   (y)  at (6, 1)  {$\hat{y}$};
%     \draw[arr] (x1) -- (h1);
%     \draw[arr] (h1) -- (y);
%   \end{tikzpicture}
 at the top of your document,
%         after \documentclass but before \begin{document}.
% =============================================================

% --- Encoding & language ---
\usepackage[T1]{fontenc}
\usepackage[utf8]{inputenc}
\usepackage[english]{babel}         % swap for [french] if needed

% --- Page geometry ---
\usepackage{geometry}
\geometry{top=2.5cm, bottom=2.5cm, left=2.5cm, right=2.5cm}

% --- Math ---
\usepackage{amsmath, amssymb}

% --- Tables ---
\usepackage{booktabs}               % \toprule, \midrule, \bottomrule
\usepackage{array}                  % column type extensions
\usepackage{tabularx}               % auto-width columns
\usepackage{multirow}               % multi-row cells

% --- Colors ---
\usepackage{xcolor}

\definecolor{mainblue}{RGB}{30, 80, 160}
\definecolor{lightblue}{RGB}{220, 235, 255}
\definecolor{accentgray}{RGB}{245, 246, 248}
\definecolor{darkgray}{RGB}{80, 80, 80}
\definecolor{solutiongreen}{RGB}{20, 120, 60}
\definecolor{lightgreen}{RGB}{220, 245, 230}
\definecolor{warnorange}{RGB}{210, 120, 20}
\definecolor{lightorange}{RGB}{255, 243, 220}

% --- Lists ---
\usepackage{enumitem}

% --- TikZ & pgfplots ---
\usepackage{tikz}
\usepackage{pgfplots}
\pgfplotsset{compat=1.18}
\usetikzlibrary{positioning, arrows.meta, calc, patterns}

% --- Boxes: tcolorbox ---
\usepackage{tcolorbox}
\tcbuselibrary{skins, breakable}

% Generic exercise box (blue border, gray background)
\tcbset{
  exercisebox/.style={
    enhanced,
    colback=accentgray,
    colframe=mainblue,
    boxrule=0.8pt,
    arc=4pt,
    left=10pt, right=10pt, top=8pt, bottom=8pt,
    fonttitle=\bfseries\color{mainblue},
  }
}

% Solution box (green, used in answer sheets)
\tcbset{
  solutionbox/.style={
    enhanced,
    colback=lightgreen,
    colframe=solutiongreen,
    boxrule=0.8pt,
    arc=4pt,
    left=10pt, right=10pt, top=8pt, bottom=8pt,
    fonttitle=\bfseries\color{solutiongreen},
  }
}

% Answer highlight box (orange, for final numeric answers)
\tcbset{
  answerbox/.style={
    enhanced,
    colback=lightorange,
    colframe=warnorange,
    boxrule=0.7pt,
    arc=4pt,
    left=8pt, right=8pt, top=5pt, bottom=5pt,
    fonttitle=\bfseries\color{warnorange},
  }
}

% Step box (light blue border, gray background, for intermediate steps)
\tcbset{
  stepbox/.style={
    enhanced,
    colback=accentgray,
    colframe=mainblue!40,
    boxrule=0.5pt,
    arc=3pt,
    left=8pt, right=8pt, top=6pt, bottom=6pt,
  }
}

% --- Inline note / warning box (mdframed) ---
\usepackage{mdframed}
\newmdenv[
  linecolor=warnorange,
  backgroundcolor=orange!5,
  roundcorner=4pt,
  innerleftmargin=8pt,
  innerrightmargin=8pt,
  innertopmargin=6pt,
  innerbottommargin=6pt,
  linewidth=0.7pt
]{notebox}

% --- Section formatting ---
\usepackage{titlesec}

% Section style: "Exercise N : Title" with a blue rule underneath
% Change the prefix via \renewcommand{\sectionprefix}{...} before \section
\newcommand{\sectionprefix}{Exercise}
\titleformat{\section}[block]
  {\large\bfseries\color{mainblue}}
  {\sectionprefix~\thesection\,:}{0.6em}{}
  [\vspace{-0.3em}\textcolor{mainblue}{\rule{\linewidth}{0.8pt}}]
\titlespacing{\section}{0pt}{1.6em}{0.8em}

% --- Header / Footer ---
\usepackage{fancyhdr}
\pagestyle{fancy}
\fancyhf{}
\renewcommand{\headrulewidth}{0.4pt}

% Customise these three commands in your document after % =============================================================
%  preamble.tex — Reusable preamble for Neural Networks course
%  Usage: \input{preamble} at the top of your document,
%         after \documentclass but before \begin{document}.
% =============================================================

% --- Encoding & language ---
\usepackage[T1]{fontenc}
\usepackage[utf8]{inputenc}
\usepackage[english]{babel}         % swap for [french] if needed

% --- Page geometry ---
\usepackage{geometry}
\geometry{top=2.5cm, bottom=2.5cm, left=2.5cm, right=2.5cm}

% --- Math ---
\usepackage{amsmath, amssymb}

% --- Tables ---
\usepackage{booktabs}               % \toprule, \midrule, \bottomrule
\usepackage{array}                  % column type extensions
\usepackage{tabularx}               % auto-width columns
\usepackage{multirow}               % multi-row cells

% --- Colors ---
\usepackage{xcolor}

\definecolor{mainblue}{RGB}{30, 80, 160}
\definecolor{lightblue}{RGB}{220, 235, 255}
\definecolor{accentgray}{RGB}{245, 246, 248}
\definecolor{darkgray}{RGB}{80, 80, 80}
\definecolor{solutiongreen}{RGB}{20, 120, 60}
\definecolor{lightgreen}{RGB}{220, 245, 230}
\definecolor{warnorange}{RGB}{210, 120, 20}
\definecolor{lightorange}{RGB}{255, 243, 220}

% --- Lists ---
\usepackage{enumitem}

% --- TikZ & pgfplots ---
\usepackage{tikz}
\usepackage{pgfplots}
\pgfplotsset{compat=1.18}
\usetikzlibrary{positioning, arrows.meta, calc, patterns}

% --- Boxes: tcolorbox ---
\usepackage{tcolorbox}
\tcbuselibrary{skins, breakable}

% Generic exercise box (blue border, gray background)
\tcbset{
  exercisebox/.style={
    enhanced,
    colback=accentgray,
    colframe=mainblue,
    boxrule=0.8pt,
    arc=4pt,
    left=10pt, right=10pt, top=8pt, bottom=8pt,
    fonttitle=\bfseries\color{mainblue},
  }
}

% Solution box (green, used in answer sheets)
\tcbset{
  solutionbox/.style={
    enhanced,
    colback=lightgreen,
    colframe=solutiongreen,
    boxrule=0.8pt,
    arc=4pt,
    left=10pt, right=10pt, top=8pt, bottom=8pt,
    fonttitle=\bfseries\color{solutiongreen},
  }
}

% Answer highlight box (orange, for final numeric answers)
\tcbset{
  answerbox/.style={
    enhanced,
    colback=lightorange,
    colframe=warnorange,
    boxrule=0.7pt,
    arc=4pt,
    left=8pt, right=8pt, top=5pt, bottom=5pt,
    fonttitle=\bfseries\color{warnorange},
  }
}

% Step box (light blue border, gray background, for intermediate steps)
\tcbset{
  stepbox/.style={
    enhanced,
    colback=accentgray,
    colframe=mainblue!40,
    boxrule=0.5pt,
    arc=3pt,
    left=8pt, right=8pt, top=6pt, bottom=6pt,
  }
}

% --- Inline note / warning box (mdframed) ---
\usepackage{mdframed}
\newmdenv[
  linecolor=warnorange,
  backgroundcolor=orange!5,
  roundcorner=4pt,
  innerleftmargin=8pt,
  innerrightmargin=8pt,
  innertopmargin=6pt,
  innerbottommargin=6pt,
  linewidth=0.7pt
]{notebox}

% --- Section formatting ---
\usepackage{titlesec}

% Section style: "Exercise N : Title" with a blue rule underneath
% Change the prefix via \renewcommand{\sectionprefix}{...} before \section
\newcommand{\sectionprefix}{Exercise}
\titleformat{\section}[block]
  {\large\bfseries\color{mainblue}}
  {\sectionprefix~\thesection\,:}{0.6em}{}
  [\vspace{-0.3em}\textcolor{mainblue}{\rule{\linewidth}{0.8pt}}]
\titlespacing{\section}{0pt}{1.6em}{0.8em}

% --- Header / Footer ---
\usepackage{fancyhdr}
\pagestyle{fancy}
\fancyhf{}
\renewcommand{\headrulewidth}{0.4pt}

% Customise these three commands in your document after \input{preamble}:
%   \renewcommand{\doccoursename}{My Course}
%   \renewcommand{\docsessionname}{Session N — Topic}
\newcommand{\doccoursename}{Neural Networks for Engineers}
\newcommand{\docsessionname}{Session}

\fancyhead[L]{\small\textcolor{darkgray}{\doccoursename}}
\fancyhead[R]{\small\textcolor{darkgray}{\docsessionname}}
\fancyfoot[C]{\small\textcolor{darkgray}{\thepage}}

% --- Title block macro ---
% Usage: \maketitleblock{Title line}{Subtitle / metadata line}{background color}
% Examples:
%   \maketitleblock{Session 4 — Exercises}{Mixed level | Neural Networks}{mainblue}
%   \maketitleblock{Session 4 — Solutions}{Perceptrons and MLPs | Neural Networks}{solutiongreen}
\newcommand{\maketitleblock}[3]{%
  \begin{tcolorbox}[
    enhanced,
    colback=#3,
    colframe=#3,
    arc=6pt, boxrule=0pt,
    left=16pt, right=16pt, top=12pt, bottom=12pt
  ]
    {\large\bfseries\color{white} #1}\\[4pt]
    {\small\color{white!85!black} #2}
  \end{tcolorbox}
  \vspace{1em}
}

% --- Convenience macros ---

% Blank answer field (inline underline for fill-in tables)
\newcommand{\acell}{\rule{1.8cm}{0.4pt}}

% Step-by-step computation display: \step{label}{math}
% Example: \step{Hidden layer}{z^{(1)} = Wx + b = \ldots}
\newcommand{\step}[2]{%
  \noindent\textbf{#1:}\quad $#2$\par\vspace{0.3em}
}

% Neuron style for TikZ diagrams (reuse with \node[neuron])
% Colors are overridable per layer — see examples below.
%
% Input neuron  : fill=lightblue,  draw=mainblue
% Hidden neuron : fill=green!15,   draw=green!50!black
% Output neuron : fill=orange!15,  draw=orange!70!black
%
% Example architecture diagram:
%
%   \begin{tikzpicture}[
%     neuron/.style={circle, minimum size=1cm, thick, font=\small\bfseries},
%     arr/.style={->, >=Stealth, thick, gray!70}
%   ]
%     \node[neuron, fill=lightblue,  draw=mainblue]          (x1) at (0, 1)  {$x_1$};
%     \node[neuron, fill=green!15,   draw=green!50!black]    (h1) at (3, 1)  {$h_1$};
%     \node[neuron, fill=orange!15,  draw=orange!70!black]   (y)  at (6, 1)  {$\hat{y}$};
%     \draw[arr] (x1) -- (h1);
%     \draw[arr] (h1) -- (y);
%   \end{tikzpicture}
:
%   \renewcommand{\doccoursename}{My Course}
%   \renewcommand{\docsessionname}{Session N — Topic}
\newcommand{\doccoursename}{Neural Networks for Engineers}
\newcommand{\docsessionname}{Session}

\fancyhead[L]{\small\textcolor{darkgray}{\doccoursename}}
\fancyhead[R]{\small\textcolor{darkgray}{\docsessionname}}
\fancyfoot[C]{\small\textcolor{darkgray}{\thepage}}

% --- Title block macro ---
% Usage: \maketitleblock{Title line}{Subtitle / metadata line}{background color}
% Examples:
%   \maketitleblock{Session 4 — Exercises}{Mixed level | Neural Networks}{mainblue}
%   \maketitleblock{Session 4 — Solutions}{Perceptrons and MLPs | Neural Networks}{solutiongreen}
\newcommand{\maketitleblock}[3]{%
  \begin{tcolorbox}[
    enhanced,
    colback=#3,
    colframe=#3,
    arc=6pt, boxrule=0pt,
    left=16pt, right=16pt, top=12pt, bottom=12pt
  ]
    {\large\bfseries\color{white} #1}\\[4pt]
    {\small\color{white!85!black} #2}
  \end{tcolorbox}
  \vspace{1em}
}

% --- Convenience macros ---

% Blank answer field (inline underline for fill-in tables)
\newcommand{\acell}{\rule{1.8cm}{0.4pt}}

% Step-by-step computation display: \step{label}{math}
% Example: \step{Hidden layer}{z^{(1)} = Wx + b = \ldots}
\newcommand{\step}[2]{%
  \noindent\textbf{#1:}\quad $#2$\par\vspace{0.3em}
}

% Neuron style for TikZ diagrams (reuse with \node[neuron])
% Colors are overridable per layer — see examples below.
%
% Input neuron  : fill=lightblue,  draw=mainblue
% Hidden neuron : fill=green!15,   draw=green!50!black
% Output neuron : fill=orange!15,  draw=orange!70!black
%
% Example architecture diagram:
%
%   \begin{tikzpicture}[
%     neuron/.style={circle, minimum size=1cm, thick, font=\small\bfseries},
%     arr/.style={->, >=Stealth, thick, gray!70}
%   ]
%     \node[neuron, fill=lightblue,  draw=mainblue]          (x1) at (0, 1)  {$x_1$};
%     \node[neuron, fill=green!15,   draw=green!50!black]    (h1) at (3, 1)  {$h_1$};
%     \node[neuron, fill=orange!15,  draw=orange!70!black]   (y)  at (6, 1)  {$\hat{y}$};
%     \draw[arr] (x1) -- (h1);
%     \draw[arr] (h1) -- (y);
%   \end{tikzpicture}
:
%   \renewcommand{\doccoursename}{My Course}
%   \renewcommand{\docsessionname}{Session N — Topic}
\newcommand{\doccoursename}{Neural Networks for Engineers}
\newcommand{\docsessionname}{Session}

\fancyhead[L]{\small\textcolor{darkgray}{\doccoursename}}
\fancyhead[R]{\small\textcolor{darkgray}{\docsessionname}}
\fancyfoot[C]{\small\textcolor{darkgray}{\thepage}}

% --- Title block macro ---
% Usage: \maketitleblock{Title line}{Subtitle / metadata line}{background color}
% Examples:
%   \maketitleblock{Session 4 — Exercises}{Mixed level | Neural Networks}{mainblue}
%   \maketitleblock{Session 4 — Solutions}{Perceptrons and MLPs | Neural Networks}{solutiongreen}
\newcommand{\maketitleblock}[3]{%
  \begin{tcolorbox}[
    enhanced,
    colback=#3,
    colframe=#3,
    arc=6pt, boxrule=0pt,
    left=16pt, right=16pt, top=12pt, bottom=12pt
  ]
    {\large\bfseries\color{white} #1}\\[4pt]
    {\small\color{white!85!black} #2}
  \end{tcolorbox}
  \vspace{1em}
}

% --- Convenience macros ---

% Blank answer field (inline underline for fill-in tables)
\newcommand{\acell}{\rule{1.8cm}{0.4pt}}

% Step-by-step computation display: \step{label}{math}
% Example: \step{Hidden layer}{z^{(1)} = Wx + b = \ldots}
\newcommand{\step}[2]{%
  \noindent\textbf{#1:}\quad $#2$\par\vspace{0.3em}
}

% Neuron style for TikZ diagrams (reuse with \node[neuron])
% Colors are overridable per layer — see examples below.
%
% Input neuron  : fill=lightblue,  draw=mainblue
% Hidden neuron : fill=green!15,   draw=green!50!black
% Output neuron : fill=orange!15,  draw=orange!70!black
%
% Example architecture diagram:
%
%   \begin{tikzpicture}[
%     neuron/.style={circle, minimum size=1cm, thick, font=\small\bfseries},
%     arr/.style={->, >=Stealth, thick, gray!70}
%   ]
%     \node[neuron, fill=lightblue,  draw=mainblue]          (x1) at (0, 1)  {$x_1$};
%     \node[neuron, fill=green!15,   draw=green!50!black]    (h1) at (3, 1)  {$h_1$};
%     \node[neuron, fill=orange!15,  draw=orange!70!black]   (y)  at (6, 1)  {$\hat{y}$};
%     \draw[arr] (x1) -- (h1);
%     \draw[arr] (h1) -- (y);
%   \end{tikzpicture}


\renewcommand{\doccoursename}{Neural Networks for Engineers}
\renewcommand{\docsessionname}{Session 6 — Quiz}

\begin{document}

\maketitleblock{Session 6 — Quiz}{Backpropagation \& Chain Rule \quad|\quad 16 questions}{mainblue}

\renewcommand{\sectionprefix}{Question}

% ------------------------------------------------------------------
\section{MSE by Hand}

A model produces the following predictions. Compute the MSE.

\begin{center}
\begin{tabular}{cccc}
\toprule
Sample & True $y$ & Predicted $\hat{y}$ & Squared Error \\
\midrule
1 & 2.0 & 2.5 & \acell \\
2 & 4.0 & 3.0 & \acell \\
3 & 1.0 & 1.5 & \acell \\
\bottomrule
\end{tabular}
\end{center}

\vspace{0.5em}
\textbf{MSE} $=$ \acell

% ------------------------------------------------------------------
\section{Why Squaring?}

Why do we square the errors in MSE rather than simply averaging them?

\vspace{5em}

% ------------------------------------------------------------------
\section{Derivative of a Loss}

Given $L(w) = (4 - 2w)^2$, compute $\dfrac{dL}{dw}$.

\vspace{5em}

% ------------------------------------------------------------------
\section{Gradient Direction}

You are at $w = 0$ on the loss curve of Question~3. The gradient is negative.
Should you \textbf{increase} or \textbf{decrease} $w$ to reduce the loss? Explain why.

\vspace{5em}

% ------------------------------------------------------------------
\section{Manual Gradient Descent}

Perform \textbf{2 steps} of gradient descent to minimize $L(w) = w^2 - 6w + 10$,
starting at $w_0 = 0$ with learning rate $\eta = 0.2$.

\begin{center}
\begin{tabular}{ccccc}
\toprule
Step & $w$ & $L(w)$ & $\dfrac{dL}{dw}$ & New $w$ \\
\midrule
0 & 0 & \acell & \acell & \acell \\[1.2em]
1 & \acell & \acell & \acell & \acell \\[1.2em]
\bottomrule
\end{tabular}
\end{center}

% ------------------------------------------------------------------
\section{Loss Landscape}

What is the difference between a \textbf{local minimum} and a \textbf{global minimum}
on a loss landscape? For which type of model is the MSE loss guaranteed to have
only one minimum?

\vspace{6em}

% ------------------------------------------------------------------
\section{Partial Derivatives}

Given $L(w_1, w_2) = (2 - w_1 - 3w_2)^2$, compute the gradient $\nabla L$
and evaluate it at $(w_1, w_2) = (1,\, 0)$.

\vspace{6em}

% ------------------------------------------------------------------
\section{SGD Variants}

Complete the table.

\begin{center}
\begin{tabularx}{\linewidth}{lccX}
\toprule
Variant & Batch size & Gradient estimate & Typical use case \\
\midrule
Full-batch GD   & $N$        & \acell & \acell \\[1em]
Mini-batch SGD  & $B \ll N$  & \acell & \acell \\[1em]
Stochastic GD   & \acell     & \acell & \acell \\[1em]
\bottomrule
\end{tabularx}
\end{center}

% ------------------------------------------------------------------
\section{Diverging Loss}

A colleague trains a linear regression model and observes that the loss
\textbf{increases} over the first few epochs. What is the most likely cause,
and how would you fix it?

\vspace{6em}

% ------------------------------------------------------------------
\section{Bridge to Neural Networks}

The notebook argues that gradient descent on linear regression and on a
multi-layer neural network follow the \textbf{same recipe}, with one key
difference. What is that difference, and what technique (covered in Session~6)
addresses it?

\vspace{6em}

% ==================================================================
% SESSION 6 — BACKPROPAGATION
% ==================================================================

% ------------------------------------------------------------------
\section{Chain Rule — Symbolic Derivative}

Let $y = \sigma(3x - 2)$, where $\sigma$ denotes the sigmoid function
and $\dfrac{d\sigma}{dz} = \sigma(z)\bigl(1 - \sigma(z)\bigr)$.

Compute $\dfrac{dy}{dx}$.

\vspace{6em}

% ------------------------------------------------------------------
\section{Chain Rule — Numerical}

Given the composite function $y = (2x + 1)^3$, evaluate
$\dfrac{dy}{dx}$ at $x = 1$.

\vspace{5em}

% ------------------------------------------------------------------
\section{Forward Pass by Hand}
\label{sec:forward-pass-by-hand}

A single neuron has weight $w = 0.8$, bias $b = -0.3$,
and uses a sigmoid activation.
Given input $x = 1.5$ and target $y = 1$:

\begin{enumerate}[label=\alph*.]
  \item Fill in the forward-pass table.
  \item Compute the MSE loss $L = (\hat{y} - y)^2$.
\end{enumerate}

\begin{center}
\begin{tabular}{cccc}
\toprule
Step & Formula & Expression & Value \\
\midrule
1 & $z = wx + b$ & $0.8 \times 1.5 - 0.3$ & \acell \\[0.8em]
2 & $\hat{y} = \sigma(z)$ & $\sigma(\;\acell\;)$ & \acell \\[0.8em]
3 & $e = \hat{y} - y$ & & \acell \\[0.8em]
4 & $L = e^2$ & & \acell \\[0.8em]
\bottomrule
\end{tabular}
\end{center}

% ------------------------------------------------------------------
\section{Backward Pass by Hand}

Continue from Question~\ref{sec:forward-pass-by-hand} (use the values you
computed).
Fill in the backward-pass table to find the gradients
$\dfrac{\partial L}{\partial w}$ and $\dfrac{\partial L}{\partial b}$.

\begin{center}
\begin{tabular}{ccc}
\toprule
Step & Gradient & Value \\
\midrule
1 & $\dfrac{\partial L}{\partial e} = 2e$ & \acell \\[1.2em]
2 & $\dfrac{\partial L}{\partial \hat{y}} = \dfrac{\partial L}{\partial e} \cdot 1$ & \acell \\[1.2em]
3 & $\dfrac{\partial L}{\partial z} = \dfrac{\partial L}{\partial \hat{y}} \cdot \sigma(z)\bigl(1-\sigma(z)\bigr)$ & \acell \\[1.2em]
4 & $\dfrac{\partial L}{\partial w} = \dfrac{\partial L}{\partial z} \cdot x$ & \acell \\[1.2em]
5 & $\dfrac{\partial L}{\partial b} = \dfrac{\partial L}{\partial z} \cdot 1$ & \acell \\[1.2em]
\bottomrule
\end{tabular}
\end{center}

After one gradient-descent step with $\eta = 1.0$, what are the updated
values of $w$ and $b$?

\vspace{4em}

% ------------------------------------------------------------------
\section{Error Signal $\delta$}

Consider a 2-layer network (2 hidden neurons, 1 output, all sigmoid).
The quantities below are known after a forward pass:

\[
  \delta^{(2)} = -0.20, \quad
  W^{(2)} = \begin{bmatrix} 0.6 \\ -0.4 \end{bmatrix},\quad
  z^{(1)} = \begin{bmatrix} 0.6 \\ -0.3 \end{bmatrix}
\]

\begin{enumerate}[label=\alph*.]
  \item Compute $(W^{(2)})^T\, \delta^{(2)}$.
  \item Compute $\sigma'(z^{(1)})$ element-wise
        (recall $\sigma'(z) = \sigma(z)(1-\sigma(z))$).
  \item Deduce $\delta^{(1)} = (W^{(2)})^T \delta^{(2)} \odot \sigma'(z^{(1)})$.
\end{enumerate}

\vspace{8em}

% ------------------------------------------------------------------
\section{Vanishing Gradients — Conceptual}

The sigmoid derivative satisfies
$\sigma'(z) \leq 0.25$ for all $z$.

\begin{enumerate}[label=\alph*.]
  \item Explain what happens to the gradient magnitude as it is
        back-propagated through many sigmoid layers.
        \vspace{4em}
  \item Name one activation function that was introduced to mitigate
        this problem, and state its derivative.
        \vspace{4em}
\end{enumerate}

\end{document}
